\subsection{Partial descent results for the henselized candidate}

The arbitrary-base fpqc descent question for
$P(R)=\alpha_p(R^h_{(t)})$ is resolved by the finite-module argument
in Theorem~\ref{A:hensel-fpqc} below.  We retain the following
independent partial descent statements, established while examining
the alternative open-relation approach.  They require neither
noetherianity nor finite presentation.

For a closed subscheme $Z\subset T$ with ideal sheaf $\mathcal I$, write
$Z^{\mathrm{oi}}$ for the union of the open subschemes $V\subset T$
on which $\mathcal I|_V=0$.  This is the \emph{schematic open interior}
of $Z$ in $T$.  It must not be confused with the topological interior
of $|Z|$: a nonzero nilpotent ideal can have support all of $T$.

\begin{lemma}[Flat factorization through a schematic open interior]
\label{A:fpqc-schematic-interior}
\status{Proved}
Let $Z\subset T$ be a closed subscheme defined by a locally finitely
generated ideal sheaf.  Every flat morphism $f:Y\to T$ which factors
through $Z$ factors through $Z^{\mathrm{oi}}$.
\end{lemma}
\begin{proof}
Fix $y\in Y$ and put $v=f(y)$.  The flat local homomorphism
$\mathcal O_{T,v}\to\mathcal O_{Y,y}$ is faithfully flat.
The image of $\mathcal I_v$ is zero because $f$ factors through $Z$.
Faithful flatness therefore gives $\mathcal I_v=0$.
Choose finitely many local generators of $\mathcal I$ near $v$.
Each generator vanishes on a neighborhood of $v$, and the intersection
of these finitely many neighborhoods is an open neighborhood on which
$\mathcal I$ is zero.  Hence $v\in Z^{\mathrm{oi}}$.
This holds for every $y$, so $f$ factors through that open subscheme.
The argument does not assert that an arbitrary flat morphism is open.
\end{proof}

\begin{proposition}[An open equality relation]
\label{A:fpqc-open-relation}
\status{Proved}
Let $A\to B$ be flat, put $X=\Spec B$, and fix $b\in B$.
Let $R_b\subset X\times_A X$ be the closed subscheme defined by
$\delta b=b\otimes1-1\otimes b$, and put $O_b=R_b^{\mathrm{oi}}$.
If the diagonal $X\to X\times_A X$ factors through $O_b$, then
$O_b\rightrightarrows X$ is an open equivalence relation.

In particular, this conclusion holds if $t\in A$ has image in
$\operatorname{Jac}(B)$ and there is a $c\in B\otimes_A B$ such that
\[
                 (1+tc)\,\delta b=0.                         \tag{A.f1}\label{eq:A-fpqc-annihilator}
\]
In that case $O_b$ also contains the entire closed subscheme
$(X\times_A X)\times_A\Spec(A/tA)$.
\end{proposition}
\begin{proof}
Symmetry follows from the involution interchanging the two factors:
it takes $\delta b$ to $-\delta b$ and preserves $O_b$.
Reflexivity is the assumed factorization of the diagonal.
For transitivity, in $X\times_A X\times_A X$ put
\[
          Y=\operatorname{pr}_{12}^{-1}(O_b)
             \cap\operatorname{pr}_{23}^{-1}(O_b).
\]
The projection $\operatorname{pr}_{13}$ is flat, being a base change
of $X\to\Spec A$, and its restriction to the open subscheme $Y$
is flat.  On $Y$ the first and second pullbacks of $b$ agree, as do
the second and third.  Thus $\operatorname{pr}_{13}|_Y$ factors
through the closed subscheme $R_b$.  Its defining ideal is principal,
so Lemma~\ref{A:fpqc-schematic-interior} makes this morphism factor
through $O_b$.  This is transitivity as a statement about schemes,
and hence about points with values in every test scheme.

For the last assertion set $s=1+tc$.  Equation~\eqref{eq:A-fpqc-annihilator} gives
$D(s)\subset O_b$.  Under the multiplication map
$\mu:B\otimes_A B\to B$, the element $s$ maps to
$1+t\mu(c)$, a unit because $t\in\operatorname{Jac}(B)$.
Consequently the diagonal factors through $D(s)$.
Finally $s$ restricts to $1$ modulo $t$, so the entire closed
$t$-fiber factors through $D(s)$ and therefore through $O_b$.
\end{proof}

\begin{proposition}[Vanishing over every reduced base]
\label{A:fpqc-reduced-base}
\status{Proved}
Let $A$ be any reduced ring and let $B$ be a flat $A$-algebra.
Suppose $b\in B$ is nilpotent and $\delta b$ vanishes on an open
neighborhood of the diagonal in $\Spec(B\otimes_A B)$.
Then $b=0$.
In particular, a nilpotent $b$ satisfying the hypotheses of the
second part of Proposition~\ref{A:fpqc-open-relation} is zero over
every reduced base $A$.
\end{proposition}
\begin{proof}
First suppose $A=k$ is a field.  For each $\mathfrak q\in\Spec B$
choose an algebraically closed extension $K$ of $\kappa(\mathfrak q)$
and the resulting homomorphism $v:B\to K$.
Base change the second factor of $\Spec(B\otimes_k B)$ by
$\Spec K\to\Spec B$.  Since $b$ is nilpotent, $v(b)=0$.
The pullback of the given open neighborhood of the diagonal is
therefore an open subscheme of $\Spec(B\otimes_k K)$ on which
$b\otimes1$ is zero.  It contains the point $w$ defined by
\[
 B\otimes_k K\longrightarrow K,
          \qquad a\otimes\lambda\longmapsto v(a)\lambda:
\]
its image in $\Spec(B\otimes_k B)$ lies on the diagonal.
The homomorphism $B\to B\otimes_k K$ is flat, and the induced local
homomorphism $B_{\mathfrak q}\to(B\otimes_k K)_w$ is faithfully
flat.  The image of $b$ in the latter ring is zero, so
$b_{\mathfrak q}=0$.  This holds for every $\mathfrak q$, proving
the assertion over a field.

For general reduced $A$, let
\[
       A'=\prod_{\mathfrak p\in\operatorname{Min}(A)}
                     \kappa(\mathfrak p).
\]
The natural homomorphism $A\to A'$ is injective: an element which
maps to zero in $\kappa(\mathfrak p)=\operatorname{Frac}(A/\mathfrak p)$
lies in $\mathfrak p$, and the intersection of the minimal primes
of a reduced ring is zero.  Flatness of $B$ gives an injection
$B\to B'=B\otimes_A A'$.
Every localization of $A'$ at a prime is a field.  Indeed, for each
$a=(a_{\mathfrak p})\in A'$ choose $d$ coordinatewise by
$d_{\mathfrak p}=a_{\mathfrak p}^{-1}$ when
$a_{\mathfrak p}\ne0$, and $d_{\mathfrak p}=0$ otherwise.
Then $a=a^2d$.  This identity persists under localization; in a
local ring with this property, $ad$ is an idempotent and every
nonzero $a$ is a unit.  Thus such a local ring is a field.
This is the usual absolutely flat reduction; compare
\cite[Lemmas 15.106.6 and 15.106.8]{StacksAbsolutelyFlatFinish}.

The nilpotence and open-neighborhood hypotheses are preserved by
arbitrary base change.  For every $\mathfrak r\in\Spec A'$, the
field case applied over $A'_{\mathfrak r}$ kills the image of $b$
in $B'\otimes_{A'}A'_{\mathfrak r}$.
An element of an $A'$-module which is zero in all localizations
at primes is zero.  Hence $b=0$ in $B'$, and the injection
$B\hookrightarrow B'$ gives $b=0$ in $B$.
No commutation of tensor product with an infinite product is used.
\end{proof}

\begin{corollary}[Reduction to a finite nilpotent ideal of the base]
\label{A:fpqc-nilpotent-base-reduction}
\status{Proved}
Let $A\to B$ be flat, with no assumption that $A$ is reduced.
If a nilpotent $b\in B$ has $\delta b=0$ on an open neighborhood
of the diagonal, then
\[
                    b\in\operatorname{Nil}(A)B.
\]
There is consequently a finitely generated nilpotent ideal
$J\subset A$ such that $b\in JB$.
\end{corollary}
\begin{proof}
Base change to $A_{\mathrm{red}}=A/\operatorname{Nil}(A)$.
The resulting algebra is flat, and the image of $b$ satisfies
the hypotheses of Proposition~\ref{A:fpqc-reduced-base}.
Its vanishing says exactly that $b\in\operatorname{Nil}(A)B$.
Write $b=\sum_{i=1}^r a_i u_i$ with $a_i\in\operatorname{Nil}(A)$
and $u_i\in B$.  Then $J=(a_1,\ldots,a_r)$ is finitely generated
and nilpotent: if $a_i^{n_i}=0$, every monomial of degree
$1+\sum_i(n_i-1)$ in the $a_i$ is zero.
\end{proof}

\paragraph{The limitation of the open-relation approach.}
In the henselian-pair reduction of the fpqc question one has a
faithfully flat map $A\to B$, henselian pairs $(A,tA)$ and $(B,tB)$,
and $b^p=0$ with \eqref{eq:A-fpqc-annihilator}.  The propositions above construct an
open equivalence relation containing the entire closed $t$-fiber
and put $b$ in the extension of a finitely generated nilpotent ideal
of $A$.  They do not show that $b$ belongs to $A$.
In particular, no effective representability or weak \emph{\'etale}
property of the fpqc quotient $X/O_b$ is asserted here.
The relation maps need not be quasi-compact or locally of finite
presentation.  Descending the coefficients of $b$ through the
nilpotent base, or constructing a suitable effective quotient
without those additional hypotheses, was the remaining step in
this approach.  Theorem~\ref{A:hensel-fpqc} resolves the descent
question by a different finite-module argument, without asserting
effectivity of this quotient.
